\documentclass[11pt]{article}
\usepackage[margin=1in]{geometry}
\usepackage{amsmath,amssymb,amsthm}
\usepackage{booktabs}
\usepackage[colorlinks=true,linkcolor=blue,citecolor=blue,urlcolor=blue]{hyperref}

\newtheorem{theorem}{Theorem}
\newtheorem{conjecture}[theorem]{Conjecture}
\newtheorem{remark}[theorem]{Remark}

\title{Disproof of TLMC Conjecture 00000000521:\\
Alternating Betti Sums of Edge Ideals}
\author{TLMC falsification submission}
\date{October 2, 2026}

\begin{document}
\maketitle

\begin{abstract}
We disprove TLMC conjecture 00000000521, which claims that for every edge
ideal $I$ the alternating Betti sums $A_j=\sum_i(-1)^i\beta_{i,j}(I)$ change
sign exactly once as $j$ crosses $\operatorname{reg}(I)$, with $|A_j|$
monotonically nondecreasing up to that point. The triangle graph
$C_3$ with $I=(xy,xz,yz)\subset k[x,y,z]$ is a counterexample: its
alternating sums are $A_0=1$, $A_1=0$, $A_2=-3$, $A_3=2$, whose signs
$+,-,+$ change twice. Even a path of two adjacent edges fails, and the
conjecture fails under either convention for reading $\beta_{i,j}(I)$.
All numbers are independently recomputed (Taylor resolution tensored with
$k$, cross-checked against Hilbert series) and formalized in axiom-free
Lean~4.
\end{abstract}

\section{The conjecture}

For a graded module let $\beta_{i,j}$ denote its graded Betti numbers and put
\[
A_j \;=\; \sum_{i\ge 0} (-1)^i\,\beta_{i,j}.
\]
Conjecture 00000000521 asserts that for every edge ideal $I$ of a simple
graph the sequence $\{A_j\}$ changes sign \emph{exactly once} as $j$ crosses
$\operatorname{reg}(I)$, and that $|A_j|$ is monotonically nondecreasing up
to that point; this was proposed as a combinatorial characterization of the
graded Lefschetz property for edge ideals.

\section{The counterexample}

Let $G=C_3$ be the triangle and let $I=I(G)=(xy,xz,yz)\subset k[x,y,z]$.
The minimal free resolution of $R/I$ is
\[
0\;\longrightarrow\; R(-3)^2
\;\longrightarrow\; R(-2)^3
\;\longrightarrow\; R
\;\longrightarrow\; R/I \;\longrightarrow\; 0,
\]
with syzygies $(z,-y,0)$, $(z,0,-x)$, $(0,y,-x)$ on the three generators and
a single syzygy among them, so the syzygy module is free of rank $2$.
Exactness is witnessed by the Hilbert series:
\[
\frac{1-3t^2+2t^3}{(1-t)^3}=\frac{(1-t)^2(1+2t)}{(1-t)^3}
=\frac{1+2t}{1-t}=1+\frac{3t}{1-t},
\]
matching $R/I\simeq k\oplus\bigoplus_{d\ge1}(kx^d\oplus ky^d\oplus kz^d)$.
Hence the Betti table of $I$ (Betti numbers of $R/I$) is
\begin{center}
\begin{tabular}{c|ccc}
 & $j=0$ & $j=2$ & $j=3$ \\ \toprule
$i=0$ & $1$ & & \\
$i=1$ & & $3$ & \\
$i=2$ & & & $2$ \\ \bottomrule
\end{tabular}
\end{center}
and therefore
\[
A_0=1,\qquad A_1=0,\qquad A_2=-3,\qquad A_3=2,\qquad A_j=0\ (j\ge4).
\]
Also $\operatorname{reg}(I)=2$: the resolution of $I$ itself is
$0\to R(-3)^2\to R(-2)^3\to I\to 0$, so $\beta_{0,2}(I)=3$,
$\beta_{1,3}(I)=2$ and $\operatorname{reg}(I)=\max(j-i)=2$.

\begin{theorem}[Falsification]
\label{thm:main}
Conjecture 00000000521 is false. For $I=(xy,xz,yz)$ the nonzero entries of
$\{A_j\}$ are $1,-3,2$, i.e.\ the signs run $+,-,+$: the sign changes
\emph{twice}, not exactly once, and $|A_0|,|A_2|,|A_3|=1,3,2$ is not
monotonically nondecreasing.
\end{theorem}

\begin{remark}[Index correction]
The original falsification-queue entry recorded the attack numbers as
$A_2=3$, $A_4=-3$, $A_6=1$, which are index-shifted. The correct numbers,
recomputed independently here (and in agreement with the verdict line), are
$A_0=1$, $A_2=-3$, $A_3=2$; the conclusion of the falsification is unchanged.
\end{remark}

\section{Boundary: how widespread the failure is}

All computations below come from the accompanying \texttt{reproduce.py}
(see Section~\ref{sec:verif}); every example passed the Hilbert-series
cross-check.

\begin{itemize}
\item The failure is minimal: already $I=(xy,yz)$ (a path of two adjacent
edges in $k[x,y,z]$, resolution $0\to R(-3)\to R(-2)^2\to R\to R/I\to0$)
gives $A=(1,0,-2,1)$: two sign changes.
\item Disjoint edges: $I=(xy,zt)\subset k[x,y,z,t]$ gives
$A=(1,0,-2,0,1)$: two sign changes. $P_4$ gives $(1,0,-3,2,0)$;
the star $K_{1,3}$ gives $(1,0,-3,3,-1)$ and $C_4$ gives
$(1,0,-4,4,-1)$ (three changes).
\item Exhaustive scan: among all $63$ labeled graphs on $4$ vertices with at
least one edge, $\mathbf{57}$ violate the ``exactly one sign change''
condition.
\item The failure is robust to the convention for $\beta_{i,j}(I)$. If
instead one resolves $I$ \emph{as a module}, then
$A_j(I)=\beta_{0,j}(R/I)-A_j(R/I)$, and: $C_3$ gives $A(I)=(0,0,3,-2)$
(one change, no violation), but $C_4$ gives $A(I)=(0,0,4,-4,1)$ and
$K_{1,3}$ gives $(0,0,3,-3,1)$: two sign changes each; a single edge gives
$A(I)=(0,0,1)$, which never changes sign. No convention rescues the
conjecture.
\end{itemize}

\section{Verification}
\label{sec:verif}

\paragraph{Independent recomputation (\texttt{reproduce.py}).}
The script builds the Taylor resolution of $R/I$, tensors it with
$k=\mathrm{GF}(10^9{+}7)$, and computes the homology dimensions per
bidegree, which yields $\beta_{i,j}(R/I)$ (Tor is computable from any free
resolution; after tensoring with $k$ only the constant entries of the
Taylor differentials survive). It then cross-checks $A_j$ against the
coefficient of $t^j$ in $(1-t)^n\cdot\operatorname{Hilb}_{R/I}(t)$ obtained
by direct monomial enumeration, for every example and every graph in the
scan. Run with \texttt{python3 reproduce.py}.

\paragraph{Lean 4 formalization (\texttt{lean4/}).}
The Betti table $\beta_{0,0}=1$, $\beta_{1,2}=3$, $\beta_{2,3}=2$, the
alternating sums $A_0=1$, $A_1=0$, $A_2=-3$, $A_3=2$, the \emph{two} sign
changes, and the failure of monotonicity of $|A_j|$ are proved by
\texttt{decide} in pure Lean~4 core (v4.33.1, no Mathlib).
\texttt{lake build \&\& lake env lean Check.lean} succeeds, and every
\texttt{\#print axioms} reports \emph{``does not depend on any axioms''};
there are no \texttt{sorry} and no additional axioms.

\end{document}
